\section{Prompt Strategy Design}
\label{sec:prompting}

Prompting is the experimental intervention. The dataset, feature set,
fingerprint, parser, ground truth, and metric suite remain fixed while the
instruction wrapper changes. The supervisor notebook originally contains a
larger strategy registry, but the validated primary results used in this
thesis contain five active strategies. Those five strategies are listed in
Table~\ref{tab:graph-prompting-strategies}.

\begin{table}[!htbp]
  \centering
  \scriptsize
  \setlength{\tabcolsep}{3pt}
  \begin{tabularx}{\textwidth}{@{}L{46mm}L{27mm}Y@{}}
    \toprule
    Strategy & Strategy role & Methodological purpose \\
    \midrule
    \shortstack[l]{\texttt{Strategy\_1}\\\texttt{Zero\_Shot\_Prompting}} & Zero-shot baseline & Uses the common base prompt without additional methodological framing. \\
    \shortstack[l]{\texttt{Strategy\_23}\\\texttt{No\_Instruction\_Control}} & Control baseline & Represents a neutral discovery instruction with no added causal-discovery policy beyond the base task. \\
    \shortstack[l]{\texttt{Strategy\_24}\\\texttt{Policy\_First\_Structured}} & Instruction structuring & Places explicit causal-discovery rules before the base task. \\
    \shortstack[l]{\texttt{Strategy\_25}\\\texttt{Least\_to\_Most}} & Prompt decomposition & Breaks graph recovery into ordered internal sub-problems from simpler physical links to more difficult accumulated-state links. \\
    \shortstack[l]{\texttt{Strategy\_26}\\\texttt{Skeleton\_of\_Thought}} & Prompt decomposition & Provides an internal answer skeleton covering energy-state edges, driving-behaviour edges, cross-group edges, and an integrity gate. \\
    \bottomrule
  \end{tabularx}
  \caption{Prompt strategies evaluated in the supervisor-provided Tesla EV results. The table uses the exact strategy names stored in the result CSV and JSON files.}
  \label{tab:graph-prompting-strategies}
\end{table}

\section{Common Base Prompt}

All evaluated strategies share a common graph-recovery scaffold. The base
prompt identifies the dataset as Tesla Full EV telematics, lists the six
features and their descriptions, supplies the dataset fingerprint, and asks
the model to identify directed feature relationships in the format
\texttt{feature\_a -> feature\_b}. The prompt explicitly requests one
relationship per line and no prose, numbering, or markdown. This output
contract is necessary because the downstream evaluator expects parseable
directed edges.

The base prompt is already a constrained analytics prompt. Therefore, the
comparison should not be interpreted as free-form language generation versus
structured generation. It is a comparison between different degrees and kinds
of instruction structure layered onto the same graph-recovery task.

\section{Baselines and Controls}

\subsection{Zero-Shot Prompting}

\strategy{Strategy_1_Zero_Shot_Prompting} applies the common base prompt
directly. It is called zero-shot because it provides no examples of correct
edges and no additional decomposition beyond the task instruction. It measures
how much graph-recovery performance the model can achieve from feature
semantics, dataset statistics, and the base output contract alone.

\subsection{No-Instruction Control}

\strategy{Strategy_23_No_Instruction_Control} is a neutral control condition.
In the active supervisor configuration, it also uses the base discovery prompt
without an added policy or decomposition layer. Its role is to provide an
additional control point for the adapted strategy set. Because it is close to
the zero-shot condition, differences between these two rows should be
interpreted conservatively and through the saved artifacts rather than treated
as evidence of a fundamentally different method.

\section{Instruction Structuring}

\strategy{Strategy_24_Policy_First_Structured} places a numbered
causal-discovery policy before the base task. The policy instructs the model
to report only direct relationships, reject correlation-only or co-movement
links, avoid bidirectional edges unless independently physically justified,
and keep the output to one ASCII-arrow edge per line.

The intended mechanism is constraint salience. By putting methodological rules
before the dataset description and task, the prompt attempts to make
directness, interventionability, and output format part of the task frame.
This is different from prompt decomposition: the policy-first prompt does not
split the task into stages, but it does make the evaluation criteria explicit
before generation.

\section{Prompt Decomposition}

Prompt decomposition strategies ask the model to organize the graph-recovery
problem into internal sub-tasks before producing the final edge list. In this
thesis, decomposition is evaluated through least-to-most prompting and
skeleton-of-thought prompting. The implementation requests only the final
edges, so the thesis evaluates observable output rather than hidden reasoning.

\subsection{Least-to-Most Prompting}

\strategy{Strategy_25_Least_to_Most} decomposes the graph-recovery problem
from simpler to more difficult sub-problems. The prompt asks the model to:
identify obvious physical edges, add edges supported by a physical law, add
edges involving accumulated state such as mileage or degradation, remove
correlation-only or mediated links, and output only the surviving final edges.

This strategy is relevant to telematics because vehicle-system relationships
are not all equally difficult. Some relationships, such as acceleration and
speed, are direct and physically familiar. Others, such as the relation
between driving behaviour and remaining energy, require integration across
energy and motion features. Least-to-most prompting tests whether an ordered
reasoning scaffold improves graph recovery compared with simpler baselines.

\subsection{Skeleton-of-Thought Prompting}

\strategy{Strategy_26_Skeleton_of_Thought} provides an internal answer
skeleton. The skeleton separates the candidate graph into energy-state edges,
driving-behaviour edges, cross-group edges, and an integrity gate that checks
for directness, correlation-only links, and invalid bidirectionality. The
model is then instructed to output only the merged final list of edges.

This strategy is more structured than least-to-most prompting because it
groups the feature graph by subsystem. In the Tesla EV task, this grouping is
natural: energy-related variables and driving-behaviour variables are
distinct, but the graph also contains cross-group relationships from speed and
acceleration to remaining energy. The strategy tests whether such an explicit
skeleton helps the model cover the graph without losing the output-format
constraint.

\section{Prompting Methods Outside the Primary Evidence}

The supervisor notebook contains commented or auxiliary strategy definitions,
including self-consistency and repeated-run machinery. Those definitions are
not treated as primary evaluated methods here because the JSON sidecars in the
validated result directory report \texttt{n\_runs = 1} for the model-strategy
cells used in the thesis. The thesis therefore does not claim measured
self-consistency improvements. It focuses on the five prompt strategies that
appear in \path{evaluation_matrix_tesla_ev_long.csv}.
